\section{Erd\H{o}s Problem \#98}

\subsection*{1) ROUND-2 OBJECTIVE}

I pursue path (A): continue the proof attempt by determining the next exact small case after Round-1. Round-1 established $h(3)=1$ and $h(4)=2$, but did not determine $h(5)$. The Round-2 target is to prove the exact value of $h(5)$.

\subsection*{2) Round-1 FOUNDATION USED}

I use the following Round-1 results without re-proving them.

\begin{itemize}
\item The exact values $h(3)=1$ and $h(4)=2$.
\item The problem remains globally unresolved: Round-1 only gave the general bounds
\[
c\,\frac{n}{\log n}\le h(n)\le n\exp(c\sqrt{\log n}).
\]
\end{itemize}

\subsection*{3) NEW INSIGHT / TOOL (ROUND-2)}

The new ingredient is a two-sided exact argument for $n=5$.

\begin{itemize}
\item The lower bound $h(5)\ge 3$ comes from a classification theorem: a $5$-point planar set with only two distances must be a regular pentagon, which is excluded by the ``no four concyclic'' hypothesis.
\item The upper bound $h(5)\le 3$ comes from an explicit $5$-point configuration in general position that determines exactly three distances.
\end{itemize}

Together these give the exact value $h(5)=3$.

\subsection*{4) ATTACK PLAN (ROUND-2)}

The gap left by Round-1 was the complete lack of exact information beyond $n=4$.

To close the next case, I need two claims:

\begin{itemize}
\item \textbf{Lower bound:} a $5$-point set in general position cannot determine only $2$ distances.
\item \textbf{Upper bound:} there exists a $5$-point set in general position determining exactly $3$ distances.
\end{itemize}

The first claim uses a precise external classification theorem; the second is fully explicit and verified by direct computation.

\subsection*{5) WORK (ROUND-2)}

\paragraph{Theorem.}
\[
h(5)=3.
\]

\paragraph{Proof of the lower bound $h(5)\ge 3$.}
I use the following theorem.

\medskip

\textbf{External theorem used.}
Among planar $5$-point sets that determine exactly two distances, the regular pentagon is the unique possibility up to similarity.

\medskip

Assume for contradiction that $h(5)=2$.
Then there exists a $5$-point set $P\subset \mathbf{R}^2$ in general position such that $|\Delta(P)|=2$.
By the external theorem, $P$ must be similar to a regular pentagon.
But all five vertices of a regular pentagon are concyclic, so $P$ violates the condition ``no four points are concyclic''.
This contradiction proves
\[
h(5)\ge 3.
\]

\paragraph{Proof of the upper bound $h(5)\le 3$.}
Consider the five points
\[
A=(0,0),\qquad
B=(1,0),\qquad
C=\left(\frac{\sqrt3}{2},\frac12\right),\qquad
D=\left(\frac12,1-\frac{\sqrt3}{2}\right),\qquad
E=(0,1).
\]
Let
\[
P:=\{A,B,C,D,E\}.
\]

I first compute all squared pairwise distances.
A direct calculation gives
\[
AB^2=AC^2=AE^2=CE^2=DE^2=1,
\]
\[
AD^2=BC^2=BD^2=CD^2=2-\sqrt3,
\]
and
\[
BE^2=2.
\]
These are all ten unordered pairs, so
\[
\Delta(P)=\left\{1,\sqrt{2-\sqrt3},\sqrt2\right\},
\]
hence
\[
|\Delta(P)|=3.
\]

It remains to check that $P$ is in general position.

\paragraph{No three points are collinear.}
For three points $X,Y,Z$, let $\Delta(X,Y,Z)$ denote twice the signed area of triangle $XYZ$.
A direct calculation yields
\[
\begin{array}{lll}
\Delta(A,B,C)=\frac12, &
\Delta(A,B,D)=1-\frac{\sqrt3}{2}, &
\Delta(A,B,E)=1, \\[6pt]
\Delta(A,C,D)=-1+\frac{\sqrt3}{2}, &
\Delta(A,C,E)=\frac{\sqrt3}{2}, &
\Delta(A,D,E)=\frac12, \\[6pt]
\Delta(B,C,D)=-\frac32+\sqrt3, &
\Delta(B,C,E)=-\frac12+\frac{\sqrt3}{2}, &
\Delta(B,D,E)=\frac12-\frac{\sqrt3}{2}, \\[6pt]
\Delta(C,D,E)=-\frac12.
\end{array}
\]
All ten values are nonzero because $\sqrt3\neq 1$ and $\sqrt3\neq 2$.
Therefore no three points of $P$ are collinear.

\paragraph{No four points are concyclic.}
For four points $X_1,\dots,X_4$ with coordinates $(x_i,y_i)$, define
\[
\Gamma(X_1,X_2,X_3,X_4):=
\det\!\begin{pmatrix}
x_1^2+y_1^2 & x_1 & y_1 & 1\\
x_2^2+y_2^2 & x_2 & y_2 & 1\\
x_3^2+y_3^2 & x_3 & y_3 & 1\\
x_4^2+y_4^2 & x_4 & y_4 & 1
\end{pmatrix}.
\]
The four points are concyclic if and only if $\Gamma=0$.
For the five $4$-subsets of $P$, direct computation gives
\[
\begin{array}{ll}
\Gamma(A,B,C,D)=\dfrac{2-\sqrt3}{2}, &
\Gamma(A,B,C,E)=\dfrac{1-\sqrt3}{2}, \\[10pt]
\Gamma(A,B,D,E)=\dfrac{1-\sqrt3}{2}, &
\Gamma(A,C,D,E)=\dfrac{\sqrt3-2}{2}, \\[10pt]
\Gamma(B,C,D,E)=\sqrt3-2.
\end{array}
\]
Again every value is nonzero.
So no four points of $P$ are concyclic.

Thus $P$ is in general position and determines exactly $3$ distances, which proves
\[
h(5)\le 3.
\]

Combining the lower and upper bounds, we conclude
\[
h(5)=3.\qquad \square
\]

\subsection*{6) ADVERSARIAL VERIFICATION}

\begin{itemize}
\item \textbf{Lower bound dependence on literature.} The only external input is the precise classification of planar $5$-point $2$-distance sets. Its hypothesis matches exactly the hypothetical case $h(5)=2$.
\item \textbf{Have all ten distances been checked?} Yes. The partition
\[
5+4+1=10
\]
of the unordered pairs matches the three displayed distance classes.
\item \textbf{Have all collinearity obstructions been checked?} Yes. There are exactly $\binom53=10$ triples, and all ten signed areas were written down explicitly.
\item \textbf{Have all cocircularity obstructions been checked?} Yes. There are exactly $\binom54=5$ four-point subsets, and all five determinants were computed explicitly.
\item \textbf{Interaction with Round-1.} This strictly extends the exact information from $h(4)=2$ to $h(5)=3$.
\end{itemize}

\subsection*{7) FINAL (EXACTLY ONE)}

\textbf{UNRESOLVED (BUT STRICTLY ADVANCED).}

Round-2 proves the exact value
\[
h(5)=3.
\]
This is a genuine new theorem beyond Round-1, though the original asymptotic question $h(n)/n\to\infty$ remains open.

\subsection*{8) COMPLETION ESTIMATE (MANDATORY)}

\textbf{COMPLETION: 40\%}

\subsection*{9) REFERENCES}

\begin{enumerate}
\item Z. Kov\'acs, G. G\'evay, and B. Kov\'acs, \emph{A note on Erd\H{o}s's mysterious remark}, arXiv:2412.05190, 2024. (Used for the uniqueness of the $5$-point planar $2$-distance configuration.)
\end{enumerate}

\bigskip
\hrule
\bigskip

